\documentclass[10pt]{beamer}
\usetheme{metropolis}
\usepackage{graphicx}
\usepackage{listings}
\usepackage{booktabs}
\usepackage{xcolor}
\definecolor{codebg}{HTML}{F5F5F5}
\definecolor{codegreen}{HTML}{2E7D32}
\definecolor{codegray}{HTML}{757575}
\definecolor{codeblue}{HTML}{1565C0}
\lstdefinestyle{rstyle}{language=R,backgroundcolor=\color{codebg},basicstyle=\ttfamily\scriptsize,
  keywordstyle=\color{codeblue}\bfseries,stringstyle=\color{codegreen},
  commentstyle=\color{codegray}\itshape,breaklines=true,frame=single,
  rulecolor=\color{codegray!30},showstringspaces=false,tabsize=2,
  morekeywords={library,ggplot,aes,geom_boxplot,geom_point,geom_jitter,geom_smooth,
    stat_summary,scale_color_manual,theme_minimal,labs,
    cooks.distance,hatvalues,rstandard,augment,influence}}
\lstdefinestyle{pythonstyle}{language=Python,backgroundcolor=\color{codebg},basicstyle=\ttfamily\scriptsize,
  keywordstyle=\color{codeblue}\bfseries,stringstyle=\color{codegreen},
  commentstyle=\color{codegray}\itshape,breaklines=true,frame=single,
  rulecolor=\color{codegray!30},showstringspaces=false,tabsize=4,
  morekeywords={import,from,as,pd,np,plt,sns,zscore,mahalanobis,
    IsolationForest,LocalOutlierFactor,EllipticEnvelope}}
\title{Outlier Detection \& Visual Diagnostics}
\subtitle{Workshop 4 --- Module 5}
\author{Prof.Asoc.\ Endri Raco}
\institute{Data Visualization for Data Scientists \\ UNYT Tirana}
\date{2025/26}
\begin{document}

\begin{frame}
  \titlepage
\end{frame}

\begin{frame}{Module Roadmap}
  \begin{enumerate}
    \item Outlier taxonomy: univariate, multivariate, contextual
    \item IQR fences and the boxplot as a detection tool
    \item Z-score vs IQR: parametric vs robust detection
    \item Mahalanobis distance: multivariate outlier detection
    \item Regression diagnostics: leverage, influence, Cook's $D$
    \item What to do with outliers: cause determines action
  \end{enumerate}
  \begin{exampleblock}{Learning Outcome}
    You will detect outliers using both robust (IQR) and parametric
    (z-score, Mahalanobis) methods, diagnose influential observations
    in regression with Cook's distance, and make principled decisions
    about whether to keep, flag, or remove each outlier.
  \end{exampleblock}
\end{frame}

\section{Outlier Taxonomy}

\begin{frame}{Three Types of Outliers}
  \begin{center}
    \includegraphics[width=0.95\textwidth]{figures_w04m05/taxonomy}
  \end{center}
  \begin{alertblock}{Pro-Tip}
    A point normal in each variable individually can still be a
    \textbf{multivariate outlier} (panel b). And a point normal
    globally can be a \textbf{contextual outlier} --- anomalous only
    in its local context (panel c). Different types require different
    detection strategies.
  \end{alertblock}
\end{frame}

\section{IQR \& Z-Score Detection}

\begin{frame}{The IQR Method: Tukey's Fences}
  \begin{center}
    \includegraphics[width=0.88\textwidth]{figures_w04m05/iqr_fences}
  \end{center}
\end{frame}

\begin{frame}{Z-Score vs IQR}
  \begin{center}
    \includegraphics[width=0.92\textwidth]{figures_w04m05/zscore_vs_iqr}
  \end{center}
  \begin{exampleblock}{Data Insight}
    Z-score assumes normality: $z = (x - \bar{x}) / s$. If the data is
    skewed, the mean and SD are pulled by outliers, making z-scores
    unreliable. IQR is \textbf{robust}: it uses medians and quartiles,
    which are resistant to extreme values. Default to IQR for EDA;
    use z-scores only for approximately normal data.
  \end{exampleblock}
\end{frame}

\begin{frame}[fragile]{Detection in Code}
  \begin{columns}[T]
    \begin{column}{0.48\textwidth}
      \textbf{R}
\begin{lstlisting}[style=rstyle]
# IQR method
q1 <- quantile(x, 0.25)
q3 <- quantile(x, 0.75)
iqr <- q3 - q1
outlier <- x < q1 - 1.5*iqr |
           x > q3 + 1.5*iqr

# Built into boxplot
ggplot(df, aes(y = Rating)) +
  geom_boxplot(outlier.colour =
    "#E53935", outlier.size = 2)

# Z-score
z <- scale(x)  # standardise
outlier_z <- abs(z) > 2
\end{lstlisting}
    \end{column}
    \begin{column}{0.48\textwidth}
      \textbf{Python}
\begin{lstlisting}[style=pythonstyle]
# IQR method
q1 = np.percentile(x, 25)
q3 = np.percentile(x, 75)
iqr = q3 - q1
outlier = (x < q1 - 1.5*iqr) | \
          (x > q3 + 1.5*iqr)

# Built into boxplot
ax.boxplot(rating,
    flierprops=dict(
        markerfacecolor="#E53935"))

# Z-score
from scipy.stats import zscore
z = zscore(x)
outlier_z = np.abs(z) > 2
\end{lstlisting}
    \end{column}
  \end{columns}
\end{frame}

\section{Multivariate Detection}

\begin{frame}{Mahalanobis Distance}
  \begin{center}
    \includegraphics[width=0.72\textwidth]{figures_w04m05/mahalanobis}
  \end{center}
  \begin{exampleblock}{Data Insight}
    Mahalanobis distance measures how far a point is from the centre
    of the data cloud, accounting for correlations. A point can be
    within normal range on each axis separately but far from the
    elliptical cloud. Threshold: $\chi^2(p, 0.95)$ where $p$ =
    number of variables.
  \end{exampleblock}
\end{frame}

\begin{frame}[fragile]{Multivariate Detection in Code}
  \begin{columns}[T]
    \begin{column}{0.48\textwidth}
      \textbf{R}
\begin{lstlisting}[style=rstyle]
# Mahalanobis (base R)
md <- mahalanobis(
  df[, c("x1", "x2")],
  colMeans(df[, c("x1", "x2")]),
  cov(df[, c("x1", "x2")]))
threshold <- qchisq(0.95, df = 2)
outlier <- md > threshold

# Robust (MCD estimator)
library(robustbase)
mcd <- covMcd(df[, c("x1","x2")])
md_robust <- mahalanobis(
  df[,c("x1","x2")],
  mcd$center, mcd$cov)
\end{lstlisting}
    \end{column}
    \begin{column}{0.48\textwidth}
      \textbf{Python}
\begin{lstlisting}[style=pythonstyle]
from scipy.spatial.distance import (
    mahalanobis)
mu = np.mean(X, axis=0)
cov = np.cov(X.T)
cov_inv = np.linalg.inv(cov)
md = [mahalanobis(xi, mu, cov_inv)
      for xi in X]
threshold = np.sqrt(
    chi2.ppf(0.95, df=2))

# Robust: EllipticEnvelope
from sklearn.covariance import (
    EllipticEnvelope)
clf = EllipticEnvelope(
    contamination=0.05)
labels = clf.fit_predict(X)
# -1 = outlier, 1 = inlier
\end{lstlisting}
    \end{column}
  \end{columns}
\end{frame}

\section{Regression Diagnostics}

\begin{frame}{Leverage vs Influence}
  \begin{center}
    \includegraphics[width=0.95\textwidth]{figures_w04m05/leverage_influence}
  \end{center}
\end{frame}

\begin{frame}{The Four Diagnostic Plots for Regression}
  \begin{center}
    \includegraphics[width=0.82\textwidth]{figures_w04m05/residual_diagnostics}
  \end{center}
\end{frame}

\begin{frame}{What Each Diagnostic Reveals}
  \centering
  \begin{tabular}{lp{5.5cm}l}
    \toprule
    \textbf{Plot} & \textbf{What It Checks} & \textbf{Bad Sign} \\
    \midrule
    (a) Resid. vs Fitted & Linearity, homoscedasticity & Funnel, curve \\
    (b) Q-Q of Residuals & Normality of errors & S-shape, heavy tails \\
    (c) Scale-Location & Constant variance (spread vs level) & Upward trend \\
    (d) Cook's Distance & Influence on fitted model & Spikes $> 4/n$ \\
    \bottomrule
  \end{tabular}
  \vspace{6pt}
  \begin{alertblock}{Pro-Tip}
    In R: \texttt{plot(lm\_model)} produces all four automatically.
    In Python: \texttt{statsmodels.graphics.plot\_regress\_exog()} or
    manual construction. Cook's $D > 4/n$ flags influential points.
    High leverage + high residual = high influence.
  \end{alertblock}
\end{frame}

\begin{frame}[fragile]{Regression Diagnostics in Code}
  \begin{columns}[T]
    \begin{column}{0.48\textwidth}
      \textbf{R}
\begin{lstlisting}[style=rstyle]
model <- lm(Rating ~ log(Reviews),
  data = apps_clean)

# All 4 plots at once
par(mfrow = c(2, 2))
plot(model)

# Or with broom + ggplot
library(broom)
diag <- augment(model)
# diag has: .fitted, .resid,
#   .hat (leverage), .cooksd

ggplot(diag, aes(x = .fitted,
    y = .resid)) +
  geom_point(alpha = 0.3) +
  geom_hline(yintercept = 0,
    linetype = "dashed") +
  geom_smooth(se = FALSE,
    color = "#E53935")
\end{lstlisting}
    \end{column}
    \begin{column}{0.48\textwidth}
      \textbf{Python}
\begin{lstlisting}[style=pythonstyle]
import statsmodels.api as sm

X = sm.add_constant(
    np.log(apps["Reviews"]+1))
model = sm.OLS(apps["Rating"],
    X).fit()

# Influence measures
infl = model.get_influence()
cooks = infl.cooks_distance[0]
leverage = infl.hat_matrix_diag

# Manual diagnostic plots
fig, axes = plt.subplots(2, 2)
axes[0,0].scatter(
    model.fittedvalues,
    model.resid, s=5, alpha=0.3)
# ... Q-Q, Scale-Loc, Cook's ...

# Quick: plot_regress_exog
from statsmodels.graphics import (
    regressionplots as rp)
rp.plot_regress_exog(model, 1)
\end{lstlisting}
    \end{column}
  \end{columns}
\end{frame}

\section{What to Do}

\begin{frame}{Cause Determines Action}
  \begin{center}
    \includegraphics[width=0.85\textwidth]{figures_w04m05/decision}
  \end{center}
\end{frame}

\begin{frame}{The ``Fit With and Without'' Protocol}
  For any influential outlier, run the analysis \textbf{twice}:
  \centering
  \begin{tabular}{lp{5cm}}
    \toprule
    \textbf{Step} & \textbf{Action} \\
    \midrule
    1 & Fit model with all data \\
    2 & Identify influential points (Cook's $D > 4/n$) \\
    3 & Refit model \textbf{without} those points \\
    4 & Compare coefficients and $R^2$ \\
    5 & Report both results \\
    \bottomrule
  \end{tabular}
  \vspace{6pt}

  If conclusions change $\Rightarrow$ the finding depends on a few
  points. This is a \textbf{fragility signal}: the result is not
  robust. Report this honestly rather than silently deleting outliers.
\end{frame}

\section{Synthesis}

\begin{frame}{Key Takeaways}
  \begin{enumerate}
    \item Three outlier types: \textbf{univariate} (extreme on one axis),
          \textbf{multivariate} (normal per-variable, extreme jointly),
          \textbf{contextual} (anomalous in local context only).
    \item \textbf{IQR fences} (Tukey) are robust to skew; \textbf{z-scores}
          assume normality. Default to IQR for EDA.
    \item \textbf{Mahalanobis distance} detects multivariate outliers by
          measuring distance from the data cloud's elliptical centre.
    \item \textbf{Leverage} $\neq$ \textbf{influence}. High leverage +
          high residual = high influence (Cook's $D$). Check with
          \texttt{plot(lm())} in R or \texttt{get\_influence()} in Python.
    \item \textbf{Never delete outliers blindly}. Determine the cause
          first: data error $\rightarrow$ fix; real extreme $\rightarrow$
          keep; influential $\rightarrow$ fit with and without.
  \end{enumerate}
\end{frame}

\begin{frame}{Essential Readings}
  \begin{itemize}
    \item Tukey, J.~W. (1977). \emph{Exploratory Data Analysis},
          Chapter 2 (Box-and-Whisker Plots, Fences).
    \item Rousseeuw, P.~J. \& van Zomeren, B.~C. (1990). ``Unmasking
          multivariate outliers.'' \emph{JASA}, 85(411).
    \item Cook, R.~D. (1977). ``Detection of influential observations
          in linear regression.'' \emph{Technometrics}, 19(1).
    \item Wilke, C.~O. (2019). \emph{Fundamentals of Data Visualization},
          Chapter 9 (Robust Statistics).
  \end{itemize}
  \vspace{6pt}
  \textbf{Next Module}: Correlation \& Association Visualization (W04-M06)
\end{frame}

\end{document}
